
\Problem{Paraulògic}

\Statement


\begin{htmlonly}
    \textbf{Llegiu aquest enunciat en PDF.}
\end{htmlonly}
    
Sí! La meva mare també està enganxada al Paraulògic...

\begin{latexonly}
    \FigureC{width=5cm}{paraulogic}
\end{latexonly}

Les regles del joc són les següents:

\begin{itemize*}
    \item Cal formar paraules amb aquestes lletres. Les lletres es poden repetir.
    \item La lletra del mig és obligatòria.
    \item Les paraules han de tenir un mínim de 3 lletres.
    \item Les paraules de 3 lletres donen 1 punt. Les de 4 lletres, 2 punts.
    \item A partir de 5 lletres, tants punts com lletres tingui la paraula.
    \item Cada joc inclou almenys un \emph{tuti} en què s’usen totes les lletres. Amb cada tuti obteniu 10 punts extra!
\end{itemize*}

Per ajudar a la meva mare feu un programa que llegeixi les set lletres de l'hexàgon i una llista de paraules d'un diccionari i escrigui les paraules puntuables del diccionari i la puntuació total.

\Input

L'entrada consisteix de paraules. La primera és de set lletres i representa la paraula de l'hexàgon, la seva primera lletra és la del mig (la roja). La resta de paraules són paraules del diccionari. Totes les lletres són en minúscules i sense accents.

\Output

La sortida són les paraules del diccionari que tenen puntació, en ordre alfabètic, seguides de cinc guions i la puntació total. Seguiu el format de l'exemple.

\clearpage
\SampleTwoCol
